% Source: https://artofproblemsolving.com/community/c3h2727233_number_theory
% Stable identity: aops-topic-2727233-q01
$S(m)>S(m-1)$ exactly when $m$ is a prime power $p^k$ (then the power of $p$ in the lcm increases), and $S(m)=S(m-1)$ otherwise. So $S(n)=S(n+4)$ exactly when none of $n+1,n+2,n+3,n+4$ is a prime power.

Checking $n=1,\dots,100$, the blocks $\{n+1,\dots,n+4\}$ free of prime powers occur exactly for
\[n=32,53,54,73,74,83,84,89,90,91,92\]
(for example $33,34,35,36$; $54,55,56,57$ and $55,56,57,58$; $90,91,92,93$ through $93,94,95,96$). So there are $\boxed{11}$ values.
